% Source pins (SHA-256): PifoStatement.lean=fe9475ebfa2c462d27cb0a99781926074e1281396f2c9e8d32d6dfca788d814e; pifo-nbe.md=5f7efd28922efc5dedad29a375f030ad6e3fe29f96ba94e77ec308f521541f14 (direct source); pifo-elegant-v4.md=356ea38146ced437e20200db04e2ce220c73f9ac9fa892b746f783b8a2d49ab8; pifo-normalization.md=2d287de7924aa30479218da3e0a9a31577ecb97c5a34395c089f7abed98f61e7 (semantic dependencies).
\section{Normalization by evaluation}
\label{sec:nbe}

We now give a second proof of \autoref{thm:main}, organized as normalization by
evaluation (NbE).  The semantic value of a scheduler is its literal flush
transformer; no topology is retained.  Semantic factor partitions then support
a canonical reification, and a common-refinement argument proves that
reification preserves every interleaved observation.

Throughout this section \(A\) is finite and carries a fixed enumeration; in the
frozen statement, \(A=\Fin(k)\).  We use the notation from
\autoref{sec:main-proof}: \(F_S:A^*\to A^*\) is the drain of a valid scheduler
\(S\), \(\restrict{w}{B}\) filters a word to \(B\subseteq A\),
\(\proj{P}(w)\) is its word of blocks for a partition \(P\), and
\(\sort_a(w)\) stably sorts occurrences by type rank \(a:A\to\Nat\) and
source arrival.  By \lemref{lem:balance},
\[
  \run(S,\flushOps(w))=\operatorname{map}(\Some,F_S(w)).
\]
Thus \(F_S(w)\) is a content-preserving permutation of the occurrences of
\(w\).  Occurrence identities need not be stored separately: equal-type
occurrences have identical paths and ranks and emerge in arrival order, so the
\(j\)-th output occurrence of a type is its \(j\)-th input occurrence.

\subsection{Drain semantics and factors}

\begin{definition}[Semantic domain and evaluation]
\label{def:nbe-semantics}
Define the realizable semantic domain by
\begin{equation}
 \mathsf{Sem}(A)
 =\{F:A^*\to A^*:F=F_S\text{ for some valid stateless scheduler }S\}.
 \label{eq:nbe-semantic-domain}
\end{equation}
Evaluation is literal draining:
\begin{equation}
  \denote{S}=F_S.
  \label{eq:nbe-evaluation}
\end{equation}
Consequently, for valid schedulers,
\[
  \denote{S}=\denote{T}
  \quad\Longleftrightarrow\quad S\flusheq T.
\]
\end{definition}

The reification phase must recover enough hierarchy from an extensional
\(F\in\mathsf{Sem}(A)\).  Its basic units are the following semantic
partitions.  For a partition \(P\), write \(P(x)\) for the unique block
containing \(x\).

\begin{definition}[Factor partition]
\label{def:nbe-factor}
A partition \(P\) of \(A\) is a \emph{factor} of \(F\in\mathsf{Sem}(A)\) when
both of the following hold.
\begin{enumerate}
\item \emph{Block locality}: for every \(B\in P\) and \(w\in A^*\),
\begin{equation}
  \restrict{F(w)}{B}=F(\restrict{w}{B}).
  \label{eq:factor-locality}
\end{equation}
\item \emph{Stable control}: some \(a:A\to\Nat\) satisfies, for every \(w\),
\begin{equation}
  \proj{P}(F(w))=\proj{P}(\sort_a(w)).
  \label{eq:factor-control}
\end{equation}
We call \(a\) a control witness.  It need not be constant within a block.
\end{enumerate}
\end{definition}

The indiscrete partition is always a factor; on the empty alphabet it is the
unique empty partition.  At an internal root, the partition into used children
is a factor.  Indeed, the root ranks give the control witness, and its stable
drain fixes the word of child invocations.  Because every push precedes every
pop, the subsequence emitted from a child is exactly the full drain of the
input subword routed there, which gives locality.  The child's global arrival
stamps may have gaps, but their order is unchanged.  A leaf will be treated as
its online-equivalent presentation with a ranked root and singleton children,
as in \lemref{lem:decomposition}.

\subsection{The factor meet}

The central semantic fact is stronger than merely amalgamating two particular
root decompositions: factors are closed under common refinement.

\begin{lemma}[Factor meet]
\label{lem:factor-meet}
Let \(F\in\mathsf{Sem}(A)\), and let \(P,Q\) be factors of \(F\).  Then
\begin{equation}
 R=P\wedge Q
   =\{B\cap C:B\in P,\ C\in Q,\ B\cap C\ne\varnothing\}
 \label{eq:factor-meet}
\end{equation}
is also a factor of \(F\).
\end{lemma}

\begin{proof}
\smallskip\noindent\emph{Locality.}
For \(D=B\cap C\in R\), filters commute, and the two factor-locality equations
give
\begin{equation}
\begin{aligned}
 \restrict{F(w)}{D}
 &=\restrict{(\restrict{F(w)}{B})}{C}
  =\restrict{F(\restrict{w}{B})}{C}\\
 &=F(\restrict{(\restrict{w}{B})}{C})
  =F(\restrict{w}{D}).
\end{aligned}
\label{eq:factor-meet-locality}
\end{equation}

\smallskip\noindent\emph{The two controls.}
Let \(a\) witness \(P\) and \(b\) witness \(Q\).  Identify each nonempty
\(R\)-cell with its pair of containing blocks.  Position by position,
\begin{equation}
 \proj{R}(F(w))
 =\operatorname{zip}\bigl(
      \proj{P}(\sort_a(w)),
      \proj{Q}(\sort_b(w))\bigr).
 \label{eq:factor-control-zip}
\end{equation}
The zip is not an arbitrary pairing: its two coordinates are the \(P\)- and
\(Q\)-projections, at the same positions, of the single permutation \(F(w)\).
It therefore has exactly the multiset of nonempty \(R\)-cells occurring in
\(w\), for every word and every choice of multiplicities.

For types in distinct \(R\)-cells, impose the exact comparison of \(a(x)\) and
\(a(y)\) when \(P\) separates \(x,y\), and the exact comparison of \(b(x)\) and
\(b(y)\) when \(Q\) separates them.  Exactness includes \(<,=,>\).  If both
partitions separate the pair, the two requirements agree.  Otherwise one of
the two arrival orders \(xy,yx\) in \autoref{eq:factor-control-zip} would pair
the diagonal coordinates into an absent off-diagonal cell, including when one
control ties the pair and the other orders it strictly.

\smallskip\noindent\emph{No finite obstruction.}
Represent a strict requirement \(x<y\) by a flagged edge \(x\to y\), and a
required tie by unflagged edges in both directions.  These constraints extend
to a weak order precisely when there is no closed directed walk containing a
flagged edge.  Suppose there is one and choose a shortest such walk.  It is a
simple cycle apart from its repeated endpoint.  A two-cycle is the pairwise
conflict just excluded.

For a longer shortest cycle, two nonconsecutive vertices in distinct
\(R\)-cells have a comparison chord.  A strict chord and the oppositely
directed cycle arc give a shorter flagged cycle; for a tie chord, choose the
arc containing the original flagged edge.  Thus every two nonconsecutive
vertices of a shortest obstruction must occupy one \(R\)-cell.  A cycle of
length at least five is then impossible: \(x_0,x_2,x_3\) would occupy one
cell, contradicting the comparison edge \(x_2\to x_3\).  In a four-cycle,
the opposite vertices form two cells.  Either \(P\) separates these cells, in
which case all four edges come from \(a\), or \(Q\) does, in which case they
all come from \(b\).  Neither weak order contains a flagged cycle.

It remains to exclude a triangle.  If either partition separates all three
pairs, one of \(a,b\) governs the triangle.  The only remaining incidence
shape, up to relabeling, is the L
\begin{equation}
 P(x)=P(y)\ne P(z),
 \qquad
 Q(y)=Q(z)\ne Q(x).
 \label{eq:factor-l-shape}
\end{equation}
Put
\begin{equation}
 (u,v,t)=
 \bigl(\cmp_b(x,y),\ \cmp_a(y,z),\ \cmp(z,x)\bigr),
 \label{eq:factor-l-comparisons}
\end{equation}
where the last comparison is common to \(a\) and \(b\).  Of the \(27\)
triples, \(13\) are compatible with a weak order.  For each of the other
\(14\), \autoref{tab:factor-l-obstructions} gives an arrival word for which the
zip in \autoref{eq:factor-control-zip} contains the missing fourth
\(P\times Q\) cell, or has the wrong cell multiplicity.  This contradicts the
all-arity content condition following \autoref{eq:factor-control-zip}.

\begin{table}[H]
\centering
\small
\begin{tabular}{@{}cc@{\qquad\qquad}cc@{}}
\toprule
\((u,v,t)\) & arrival word & \((u,v,t)\) & arrival word\\
\midrule
\((<,<,<)\) & \(xyz\) & \((>,>,>)\) & \(xyz\)\\
\((<,<,=)\) & \(yzx\) & \((>,>,=)\) & \(xyz\)\\
\((<,=,<)\) & \(xyz\) & \((>,=,>)\) & \(xzy\)\\
\((<,=,=)\) & \(yzx\) & \((>,=,=)\) & \(xzy\)\\
\((=,<,<)\) & \(xyz\) & \((=,>,>)\) & \(yxz\)\\
\((=,<,=)\) & \(zxy\) & \((=,>,=)\) & \(yxz\)\\
\((=,=,<)\) & \(xyz\) & \((=,=,>)\) & \(zyx\)\\
\bottomrule
\end{tabular}
\caption{The fourteen inconsistent L-shaped controls.  Changing the arrival
word is essential in the rows involving FIFO ties.}
\label{tab:factor-l-obstructions}
\end{table}

There is no flagged cycle.  Collapse the tie strongly connected components
and topologically order the resulting strict DAG; natural ranks for this order
give a weak priority \(g\) satisfying every imposed comparison.  Exact
cross-block comparisons determine a colored stable sort, so
\begin{equation}
 \proj{P}(\sort_g(w))=\proj{P}(\sort_a(w)),
 \qquad
 \proj{Q}(\sort_g(w))=\proj{Q}(\sort_b(w)).
 \label{eq:factor-amalgamated-controls}
\end{equation}
Zipping these equalities and using \autoref{eq:factor-control-zip} yields
\[
  \proj{R}(F(w))=\proj{R}(\sort_g(w)),
\]
which is stable control for \(R\).  Together with
\autoref{eq:factor-meet-locality}, this proves that \(R\) is a factor.
\end{proof}

The all-word condition in this proof cannot be replaced by the collection of
two-type restrictions.

\begin{example}[Pairwise flush data is insufficient]
\label{ex:pairwise-insufficient}
On types \(A,B,C\), consider the two two-leaf schedulers
\begin{equation}
\begin{array}{@{}c|c|c@{}}
 & \text{root routing and ranks} & \text{ranks in nonsingleton child}\\
\hline
 S_1 & \{A,B\}\mid\{C\},\quad A=1,\ B=C=2
     & A=B\\
 S_2 & \{A,C\}\mid\{B\},\quad A=B=C
     & A<C.
\end{array}
\label{eq:pairwise-witness-schedulers}
\end{equation}
All leaf ranks of \(S_1\) tie; the singleton leaf ranks are immaterial.  On
every word over each displayed pair, repetitions included, their drains agree:
\begin{equation}
 \{A,B\}:\mathrm{FIFO},
 \qquad
 \{A,C\}:\text{\(A\) stably before \(C\)},
 \qquad
 \{B,C\}:\mathrm{FIFO}.
 \label{eq:pairwise-witness-restrictions}
\end{equation}
Nevertheless the ternary word distinguishes them:
\begin{equation}
  F_{S_1}(CBA)=BCA,
  \qquad
  F_{S_2}(CBA)=ABC.
  \label{eq:pairwise-witness-ternary}
\end{equation}
Thus complete pairwise flush behavior does not determine a flush transformer;
the arbitrary-multiplicity content constraint in the factor-meet proof is
essential.
\end{example}

\subsection{Reification and adequacy}

There are finitely many partitions of finite \(A\).  The indiscrete partition
is a factor, and repeated application of \lemref{lem:factor-meet} shows that
the meet of all factors is a factor.  Call this finest factor \(P_F\).  When
\(A\ne\varnothing\), among the control witnesses whose ranks have been
replaced by consecutive rank classes, choose the lexicographically least
compressed weak-order witness
\[
 g_F:A\to\{0,\ldots,m-1\}.
\]
Here \(m\) is its number of nonempty rank classes.  This choice is from a
finite set and depends only on \(F\).  On the empty alphabet take the unique
empty function.

If \(|A|>1\), \(P_F\) is proper.  To see this, take any scheduler realizing
\(F\), discard unused children, and contract roots with only one used child.
A branching root supplies a proper root factor.  If the first live node is a
leaf, then \(F=\sort_a\) for its leaf ranks \(a\), and the discrete singleton
partition is a proper factor.  Since \(P_F\) refines every factor, it is proper
as well.

We may therefore define \(\mathsf{reify}_A(F)\) by recursion on \(|A|\).
For \(|A|\le1\), use a rank-zero leaf, with the vacuous assignment when \(A\)
is empty.  Otherwise, make a root with one canonically ordered child for every
\(B\in P_F\), and give type \(x\) root rank \(g_F(x)\).  The child for \(B\)
is
\[
  \mathsf{reify}_B(F|_B),
  \qquad
  F|_B:B^*\to B^*,
  \qquad
  F|_B(v)=F(v).
\]
Factor locality and content preservation make the displayed restriction
well typed.  It is realizable by restricting any witness for \(F\) to types
in \(B\), and the recursive call is smaller because \(P_F\) is proper.  The
root annotation followed by the recursively produced path is well formed, so
the result is a valid scheduler.

\begin{lemma}[Increasing arrival renaming]
\label{lem:nbe-arrival-renaming}
In a queue or tree computation, replace the global arrival stamps of source
pushes by any strictly increasing relabeling, using the same new stamp at every
node visited by one push.  Every emitted type is unchanged.
\end{lemma}

\begin{proof}
Ranks are unchanged, and a strictly increasing relabeling preserves every
lexicographic \((\rank,\arr)\) comparison.  Structural induction through the
queue pop and recursive tree pop therefore preserves each selection and
emitted type.  This is the structural fact needed for embedded child states.
\end{proof}

\begin{theorem}[Reification adequacy]
\label{thm:nbe-adequacy}
For every \(F\in\mathsf{Sem}(A)\),
\begin{equation}
  F_{\mathsf{reify}_A(F)}=F.
  \label{eq:nbe-adequacy}
\end{equation}
\end{theorem}

\begin{proof}
Proceed by induction on \(|A|\).  The base cases are immediate.  For a word
\(w\), the reified root invokes blocks in the order
\begin{equation}
 \proj{P_F}(\sort_{g_F}(w))=\proj{P_F}(F(w)).
 \label{eq:nbe-root-skeleton}
\end{equation}
For \(B\in P_F\), let \(h_B(j)\) be the one-based global position in \(w\) of
the \(j\)-th letter of \(\restrict{w}{B}\).  After all pushes, the embedded
child state is its standalone state with local arrival \(j\) replaced by the
strictly increasing \(h_B(j)\).  Apply
\lemref{lem:nbe-arrival-renaming} before the induction hypothesis.  The child
then supplies the output subsequence
\[
 F(\restrict{w}{B})=\restrict{F(w)}{B},
\]
where the equality is factor locality.  The block skeleton in
\autoref{eq:nbe-root-skeleton}, together with all of these block subsequences,
uniquely reconstructs \(F(w)\).  This proves \autoref{eq:nbe-adequacy}.
\end{proof}

\subsection{From batch control to online control}

One queue-level fact upgrades factor control to arbitrary interleavings.  It
is the semantic form of \lemref{lem:colored}.

\begin{lemma}[Batch skeletons determine online colors]
\label{lem:nbe-colored}
Let \(c:A\to C\) be a coloring and let \(a,g:A\to\Nat\) be weak priorities.
Extend \(c\) pointwise to words.
If
\begin{equation}
 c(\sort_a(w))=c(\sort_g(w))
 \label{eq:nbe-colored-batch}
\end{equation}
for every word \(w\), then two priority queues ranked by \(a\) and \(g\),
driven from empty by the same arbitrary push/pop word, emit the same color on
every pop, including \(\None\) when empty.
\end{lemma}

\begin{proof}
The two two-letter arrival orders show that
\autoref{eq:nbe-colored-batch} is equivalent to equality of the exact comparison
on every pair of differently colored types.  Normalize either weak order into
\emph{observable bands}.  Retain a rank level containing multiple colors as a
mixed FIFO band, whose abstract state is its FIFO color word.  Collapse every
maximal consecutive run of levels containing only one color into a
monochromatic band, whose state is only its pending count.

The exact comparisons with types of other colors give each type an external
signature.  These signatures, together with forced cross-color ties,
reconstruct the same ordered bands for \(a\) and \(g\); a boundary invisible
to all cross-color comparisons can occur only inside one monochromatic band.
A push therefore increments the same count or appends the same color to a
mixed FIFO word.  A pop uses the first nonempty band and performs the same
abstract transition.  If all bands are empty, both queues emit \(\None\).

The concrete queues may remove different types of one color only within a
monochromatic band having one external comparison profile.  If the band ties
an external color, it is mixed instead, and the common global arrival order
selects the same oldest occurrence.  Thus the abstract coupling covers both
hidden same-color identities and FIFO ties.
\end{proof}

\subsection{Soundness through a common semantic refinement}

We prove a statement stronger than the reification equation: any two valid
schedulers with one semantic value have identical online behavior.  This is
an alternative route to \autoref{thm:main}.

\begin{theorem}[Common-refinement soundness]
\label{thm:nbe-common-refinement}
If valid stateless schedulers \(S,T\) over finite \(A\) satisfy
\(F_S=F_T\), then \(S\intereq T\).
\end{theorem}

\begin{proof}
Use strong induction on \(|A|\); the cases \(|A|\le1\) are immediate.  First
apply two online-preserving presentation changes.  A node with one used child
can be contracted because its queue only counts child calls.  A leaf with
type ranks \(a\) can be expanded to an \(a\)-ranked root routing to singleton
children: root and child select the same occurrence, including on FIFO ties.
Thus both schedulers may be presented with proper used-child root factors.

Let the common flush function be \(F\), and write those root factors as
\((P,a)\) for \(S\) and \((Q,b)\) for \(T\).  By
\lemref{lem:factor-meet}, \(R=P\wedge Q\) has a control witness \(g\).  For
each \(D\in R\), choose one scheduler \(U_D\) realizing \(F|_D\), for example
\(\mathsf{reify}_D(F|_D)\).

Fix \(B\in P\).  Build \(V_B\) with a \(g|_B\)-ranked root whose children are
the cells \(D\subseteq B\), and put \(U_D\) below the corresponding child.
The factor equations give
\begin{equation}
  F_{V_B}=F|_B.
  \label{eq:nbe-block-refinement}
\end{equation}
Indeed, the \(R\)-control equation supplies the output cell skeleton, and
factor locality supplies the output subsequence in every cell; together they
reconstruct the word.  The actual \(B\)-child of \(S\) has the same flush
function \(F|_B\).

To use the induction hypothesis in the frozen statement's finite-alphabet
presentation, reindex the subtype \(B\) by a bijection
\(\Fin(|B|)\cong B\).  Runs commute with this transport because queue
selection inspects ranks and arrivals, never packet values.  Since \(P\) is
proper, \(|B|<|A|\), so the induction hypothesis makes the actual child and
\(V_B\) interleaved-equivalent.  Substitute every such child under the
unchanged root.  Corresponding children receive the same local operation
words.  Their global arrival stamps form a strictly increasing, possibly
gapped subsequence; apply \lemref{lem:nbe-arrival-renaming} before the
induction hypothesis.  The contextual substitutions therefore preserve the
complete observation trace.

The resulting presentation of \(S\) has an outer \(a\)-queue routing to
\(P\), followed inside every block by a \(g\)-queue routing to \(R\).  Since
\(P\) coarsens \(R\), factor control gives, for every \(w\),
\begin{equation}
 \proj{P}(\sort_a(w))
 =\proj{P}(F(w))
 =\proj{P}(\sort_g(w)).
 \label{eq:nbe-root-swap}
\end{equation}
By \lemref{lem:nbe-colored}, the outer queue can be replaced, in the system
built from empty, by a \(g\)-queue without changing the stream of \(P\)-child
calls.  This constructs a new coupled run; it does not mutate a live queue
while retaining state ordered by incompatible ranks.

The consecutive \(g\)-levels now collapse occurrence-for-occurrence.  Tag
root references by global arrival and maintain that the inner queue for \(B\)
is the \(B\)-projection of the outer queue.  A push preserves this invariant.
If the outer queue removes the unique minimum
\((g(x),\arr)\) lying in \(B\), that occurrence is also the unique minimum of
the projected inner queue.  Removing both levels therefore gives one
\(g\)-ranked root routing directly to \(R\), with the fixed children \(U_D\).

Repeating the construction from \(T\), using \(Q,b\), produces the identical
one-level scheduler.  Hence \(S\intereq T\).  Every recursive call was on a
proper root block; arrival renaming was a structural queue fact, so there is
no circular appeal to the theorem.  Finally, validity makes recursive child
pops succeed whenever a root reference is present.  Root emptiness depends
only on the preceding pushes and successful pops, so excess pops emit
\(\None\) on both sides.
\end{proof}

Combining adequacy with common-refinement soundness gives the NbE equation
\begin{equation}
 S\intereq
 \mathsf{reify}_A(\denote{S}).
 \label{eq:nbe-soundness}
\end{equation}
If \(S\flusheq T\), their semantic values are equal, so the deterministic
choices above give the same canonical reification.  Applying
\autoref{eq:nbe-soundness} on both sides yields \(S\intereq T\), proving
\autoref{thm:main} by this second route.  This is normalization by evaluation,
not a scheduler rewrite system: evaluation forgets topology, factors are
properties of the denoted function, and soundness joins two presentations
through a semantic common refinement.

\subsection{Scope}

\paragraph{No decision procedure from this construction.}
The semantic object in \autoref{eq:nbe-semantic-domain} is an extensional function
on all words.  Testing factor locality and stable control, finding all factors,
and selecting \(P_F\) quantify over infinitely many words.  The reification
above is therefore a classical mathematical map, not an algorithm supplied by
this proof.  For fixed \(k>0\), pruning unused children, erasing unary used
nodes, and compressing ranks leaves finitely many reduced presentations, with
at most \(k-1\) internal nodes because every internal node splits a nonempty
support into proper supports.  For \(k=0\), the unused root collapses to a leaf.
This finiteness supplies neither an effective factor test nor a computable
normalizer or bound on a distinguishing word.  Searching words enumerates
non-equivalence, but this argument does not terminate on equivalent
schedulers.  Thus this NbE construction by itself proves no algorithmic
decidability result for \(\intereq\); the finite structural procedure in
\autoref{sec:deciding-equivalence} uses the separate pair-and-triple invariant.

\paragraph{Infinite alphabets by finite support.}
The implication itself extends propositionally to an arbitrary packet-type
alphabet \(\alpha\).  Given one finite interleaved operation word, let \(B\)
be the finite set of types it pushes and enumerate \(B\) as \(\Fin(n)\), using
a classical enumeration of this finite image (or decidable equality on the
support), not a global enumeration of \(\alpha\).
Restrict both fixed assignment maps to \(B\).  Validity is preserved, and
flush equivalence over \(\alpha\) implies flush equivalence of the restricted
schedulers.  The finite theorem gives equality on the restricted operation
word; every emitted type lies in \(B\), so transporting back gives equality
over \(\alpha\).  This is a finite-support corollary.  It constructs neither
one global canonical scheduler on an infinite alphabet nor a decision
procedure there.
